\section{Union, Intersection, Difference}
The same way numbers can be combined with operations such as addition, multiplication, subtraction, and division, we have similar operations for sets too. 
\begin{defn}
    Let $A,B$ be sets then: 
    \begin{enumerate}
        \item The \emph{union} of $A,B$ is $A\cup B=\{ x:x\in A\lor x\in B \}$
        \item The \emph{intersection} of $A,B$ is $A\cap B=\{ x:x\in A\land x\in B \}$
        \item The \emph{difference} of $A,B$ is $A-B=\{ x:x\in A\land x\notin B \}$ 
    \end{enumerate}
\end{defn}

Let us consider the following examples: 
\begin{example}
    Let $A=\{ a,b,c,d,e \},B=\{ d,e,f \},C=\{ 1,2,3 \}$: 
    \begin{enumerate}
        \item $A\cup B=\{ a,b,c,d,e,f \}$
        \item $A\cap B=\{ d,e \}$
        \item $A-B=\{ a,b,c \}$
        \item $B-A=\{ f \}$
        \item $(A-B)\cup(B-A)=\{ a,b,c \}\cup\{ f \}=\{ a,b,c,f \}$
        \item $A\cup C=\{ a,b,c,d,e,1,2,3 \}$
        \item $A\cap C=\varnothing$
        \item $A-C=\{ a,b,c,d,e \}$
        \item $(A\cap C)\cup(A-C)=\varnothing \cup \{ a,b,c,d,e \}=\{ a,b,c,d,e \}$
        \item $(A\cap B)\times B=\{ d,e \}\times \{ d,e,f \}=\{ (d,d),(d,e),(d,f),(e,d),(e,e),(e,f) \}$
        \item $(A\times C)\cap(B\times C)=\{(d,1),(d,2),(d,3),(e,1),(e,2),(e,3)\}$.
    \end{enumerate}
\end{example}

\emph{Important:} It is always true that $X\cap Y=Y\cap X, X\cup Y=Y\cup X$. In other words, these are \emph{commutatitve operations.} However, in general, $X-Y\ne Y-X$.

\subsection{Exercises From Section 5}
\begin{exercise}\label{ex:5.1}
    Suppose $A=\{ 4,3,6,7,1,9 \},B=\{ 5,6,8,4 \},C=\{ 5,8,4 \}$. Find: 
    \begin{enumerate}
        \item $A\cup B$
        \item $A\cap B$
        \item $A-B$
        \item $A-C$
        \item $B-A$
        \item $A\cap C$
        \item $B\cap C$
        \item $B\cup C$
        \item $C-B$
    \end{enumerate}
\end{exercise}\hyperref[sol:5.1]{[Go to Solution]}

\begin{exercise}\label{ex:5.2}
    Suppose $A=\{ 0,2,4,6,8 \},B=\{ 1,3,5,7 \},C=\{ 2,8,4 \}$. Use the operations from exercise~\ref{ex:5.1}\footnote{I believe the subpoints in the original text are slightly different to the subpoints of~\ref{ex:5.1}, but if you can do the above, you can certainly do those from the book too.}
\end{exercise}\hyperref[sol:5.2]{[Go to Solution]}